\documentclass[aspectratio=169]{beamer}
\usepackage{xcolor}
\usepackage{tikz}
\usetikzlibrary{patterns}
\usepackage{amsmath, amssymb, bm, amsfonts}
\usepackage{multicol}

% settings for the beamer theme
\usetheme[progressbar= frametitle]{metropolis}
%\setbeamertemplate{frame numbering}[fraction]
\setbeamertemplate{footline}{}
\useoutertheme{metropolis}
\useinnertheme{metropolis}
\usecolortheme{spruce}
\setbeamercolor{background canvas}{bg=white}
\definecolor{mygreen}{rgb}{.125, .5, .125}
\usecolortheme[named=mygreen]{structure}

% setting for the title
%\title[Short Title]{Beamer-Krumel}
\title{\Huge Vocabulario de Conjuntos}
%\subtitle{Subtitle Here}
\usepackage{graphicx} % Required for inserting images
%\author{Kenyer Aguiar}
%\institute{Institute}
\institute{\large\textbf{Objetivo de la clase:} \\[6pt]
Practicar vocabulario básico de Teoría de Conjuntos}
\date{}

% No tocar permite que align y \pause funcionen bien
\mode<presentation> { \setbeamercovered{transparent=0} }
\makeatletter
\def\pdftex@driver{pdftex.def}
\ifx\Gin@driver\pdftex@driver
  \def\pgfsys@color@unstacked#1{%
    \pdfliteral{\csname\string\color@#1\endcsname}%
  }
\fi
\makeatother
%------------------------------



% Definir color gris claro
\definecolor{grisclaro}{RGB}{235,235,235}

% Entorno personalizado sin título
\newenvironment{mibloque}
{
\begin{center}
\begin{beamercolorbox}[
    rounded=true,
    shadow=false,
    sep=1em,
    wd=0.8\textwidth
]{}
\setbeamercolor{background canvas}{bg=grisclaro}
\color{black}
\setbeamercolor{block body}{bg=grisclaro}
\begin{beamercolorbox}[rounded=false,sep=0.3em,wd=\textwidth]{block body}
}
{
\end{beamercolorbox}
\end{beamercolorbox}
\end{center}
}

%------------------------------


\newcommand{\cuadro}[1]{
\tikz[baseline=(m.base)]{
        \node[draw=mygreen, text=mygreen, inner sep=2pt, thick, anchor=base, fill = mygreen!10] (m) {#1}
        ;}\quad}



\newcommand{\ejemplo}{%
  \tikz[baseline=(m.base)]{
    \node[
      draw=mygreen,
      text=mygreen,
      inner sep=2pt,
      thick,
      anchor=base,
      fill=mygreen!10
    ] (m) {\textbf{Ejemplo}};
  }%
\,\,}

\newcommand{\sol}{%
  \tikz[baseline=(m.base)]{
    \node[
      draw=mygreen,
      text=mygreen,
      inner sep=2pt,
      thick,
      anchor=base,
      fill=mygreen!10
    ] (m) {\textbf{Solución:}};
  }%
\,\,}
        
\begin{document}
% filling for blocks
\metroset{block=fill}

\begin{frame}
\titlepage  
\end{frame}

%--------------------------------
\begin{frame}[c]{}
\begin{block}{Pregunta 1}
\large
¿Cómo se define formalmente la relación de pertenencia en teoría de conjuntos?

\begin{enumerate}[A)]
    \item Es la relación que existe entre dos conjuntos cuando uno está dentro del otro.
    \item Es la operación que combina todos los elementos de dos colecciones distintas.
    \item Es el número total de elementos distintos que posee una agrupación.
    \item Es el vínculo que se establece exclusivamente entre un objeto (elemento) y el conjunto que lo contiene.
\end{enumerate}
\end{block}
\end{frame}
%--------------------------------
\begin{frame}
\begin{block}{Pregunta 2}
Un conjunto que no contiene ningún elemento se denomina vacío. ¿Cuál es la propiedad fundamental de su cardinalidad?

\begin{enumerate}[A)]
    \item Su cardinal es indefinido.
    \item Su cardinal es exactamente igual a 0.
    \item Su cardinal es infinito pues no tiene límites.
    \item Su cardinal es 1, representando la existencia del conjunto mismo.
\end{enumerate}
\end{block}
\end{frame}

\begin{frame}
\begin{block}{Pregunta 3}
¿Qué caracteriza a dos conjuntos denominados \textbf{disjuntos}?

\begin{enumerate}[A)]
    \item Su intersección es el conjunto vacío.
    \item La unión de ambos da como resultado el conjunto universo.
    \item Uno de ellos es subconjunto del otro.
    \item Tienen exactamente la misma cantidad de elementos.
\end{enumerate}
\end{block}
\end{frame}

\begin{frame}
\begin{block}{Pregunta 4}
Si un conjunto \(A\) tiene una cantidad de elementos que se puede contar y el proceso de conteo termina, el conjunto se clasifica como:

\begin{enumerate}[A)]
    \item Potencia
    \item Unitario
    \item Finito
    \item Universo
\end{enumerate}
\end{block}

\end{frame}

\begin{frame}
\begin{block}{Pregunta 5}
Dada la definición de \textbf{Subconjunto}, ¿cuál de las siguientes afirmaciones es siempre verdadera?

\begin{enumerate}[A)]
    \item El conjunto vacío no puede ser subconjunto de un conjunto finito.
    \item Todo conjunto es subconjunto de sí mismo.
    \item Dos conjuntos disjuntos son subconjuntos entre sí.
    \item Un subconjunto siempre debe tener menos elementos que el conjunto original.
\end{enumerate}


\end{block}
\end{frame}


\begin{frame}
\begin{block}{Pregunta 6}
¿A qué se refiere el término \textbf{Conjunto Universo} en un problema determinado?

\begin{enumerate}[A)]
    \item Al conjunto de todos los números reales existentes.
    \item Al conjunto que resulta de la intersección de todos los conjuntos dados.
    \item A la unión de todos los conjuntos infinitos.
    \item Al conjunto que contiene todos los elementos relevantes para ese contexto o estudio particular.
\end{enumerate}
\end{block}
\end{frame}


\begin{frame}
\begin{block}{Pregunta 7}
El \textbf{Complemento} de un conjunto \(A\), denotado como \(A^{c}\), se define como:

\begin{enumerate}[A)]
    \item La diferencia entre el conjunto \(A\) y el conjunto vacío.
    \item El conjunto que contiene solo los elementos negativos de \(A\).
    \item El conjunto de todos los elementos que pertenecen al universo pero no pertenecen a \(A\).
    \item El conjunto de elementos que \(A\) comparte con el universo.
\end{enumerate}
\end{block}
\end{frame}

\begin{frame}
\begin{block}{Pregunta 8}
Si el cardinal de un conjunto \(A\) es \(n\), ¿cuántos elementos tiene su \textbf{Conjunto Potencia} \(P(A)\)?

\begin{enumerate}[A.]
    \item \(2^{n}\)
    \item \(n^{2}\)
    \item \(2n\)
    \item \(n-2\) 
\end{enumerate}


\end{block}
\end{frame}


\begin{frame}
\begin{block}{Pregunta 9}
El \textbf{Producto Cartesiano} \(A \times B\) es un conjunto cuyos elementos son:

\begin{enumerate}[A.]
    \item La multiplicación de los valores numéricos de los elementos de \(A\) y \(B\).
    \item La unión de los conjuntos \(A\) y \(B\) dispuesta en una línea.
    \item Conjuntos que contienen un elemento de \(A\) y uno de \(B\) sin orden específico.
    \item Pares ordenados \((a,b)\) donde \(a \in A\) y \(b \in B\).
\end{enumerate}
\end{block}
\end{frame}


\begin{frame}
\begin{block}{Pregunta 2}

\end{block}
\end{frame}\begin{frame}
\begin{block}{Pregunta 2}

\end{block}
\end{frame}




\begin{frame}
\begin{block}{Pregunta 2}

\end{block}
\end{frame}




\begin{frame}
\begin{block}{Pregunta 2}

\end{block}
\end{frame}




\begin{frame}
\begin{block}{Pregunta 2}

\end{block}
\end{frame}\begin{frame}
\begin{block}{Pregunta 2}

\end{block}
\end{frame}




\begin{frame}
\begin{block}{Pregunta 2}

\end{block}
\end{frame}



\begin{frame}
\begin{block}{Pregunta 2}

\end{block}
\end{frame}



\begin{frame}
\begin{block}{Pregunta 2}

\end{block}
\end{frame}



\begin{frame}
\begin{block}{Pregunta 2}

\end{block}
\end{frame}



\begin{frame}
\begin{block}{Pregunta 2}

\end{block}
\end{frame}



\begin{frame}
\begin{block}{Pregunta 2}

\end{block}
\end{frame}




\begin{frame}
\begin{block}{Pregunta 2}

\end{block}
\end{frame}\begin{frame}
\begin{block}{Pregunta 2}

\end{block}
\end{frame}



\begin{frame}
\begin{block}{Pregunta 2}

\end{block}
\end{frame}




\begin{frame}
\begin{block}{Pregunta 2}

\end{block}
\end{frame}




\begin{frame}
\begin{block}{Pregunta 2}

\end{block}
\end{frame}



\begin{frame}
\begin{block}{Pregunta 2}

\end{block}
\end{frame}



\begin{frame}
\begin{block}{Pregunta 2}

\end{block}
\end{frame}



\begin{frame}
\begin{block}{Pregunta 2}

\end{block}
\end{frame}




\begin{frame}
\begin{block}{Pregunta 2}

\end{block}
\end{frame}



\begin{frame}
\begin{block}{Pregunta 2}

\end{block}
\end{frame}\begin{frame}
\begin{block}{Pregunta 2}

\end{block}
\end{frame}



\begin{frame}
\begin{block}{Pregunta 2}

\end{block}
\end{frame}\begin{frame}
\begin{block}{Pregunta 2}

\end{block}
\end{frame}


\begin{frame}
\begin{block}{Pregunta 2}

\end{block}
\end{frame}


\begin{frame}
\begin{block}{Pregunta 2}

\end{block}
\end{frame}



\begin{frame}
\begin{block}{Pregunta 2}

\end{block}
\end{frame}



\begin{frame}
\begin{block}{Pregunta 2}

\end{block}
\end{frame}



\begin{frame}
\begin{block}{Pregunta 2}

\end{block}
\end{frame}


\begin{frame}
\begin{block}{Pregunta 2}

\end{block}
\end{frame}


\end{document}

